\subsection{THETA3-C26: an integral unit on a two-copy isotropic CM direction}\label{C26}

\paragraph{Statement and normalization.}
Use the elliptic CM hypotheses of C19: $p\geq5$ is inert in the
imaginary quadratic CM field $K$, the CM order $\mathcal O$ is
maximal at $p$, and $E$ has a good model over a complete discrete
valuation ring $R$ finite over $\Z_p$.  All CM endomorphisms are
defined, and their action on the chosen integral invariant
differential is the fixed embedding $K\hookrightarrow\operatorname{Frac}(R)$.
Let $t$ be an integral parameter, let $\lambda(t)=t+\cdots$
be its logarithm for this differential, and let $D\lambda=1$.
Use the actual reduced theta normalized by $\theta'(0)=1$ and put
\[
 q(t)=\frac{\theta(\lambda(t))}{t}\in1+t\C_p[[t]].
\]
No integrality of $q$ is assumed.

For every $b\in\mathcal O_{K_p}^{\times}$ the completed CM
action of C19 satisfies
\begin{equation}\label{C26-unit-isogeny}
 U_b(t):=\frac{q([b](t))}{q(t)^{b\bar b}}
           \in1+tR[[t]].
\end{equation}
The exponent on the right is a $p$-adic integer and is interpreted
coefficientwise by the binomial series at the constant term one.
In particular, if $b\bar b=-1$, then
\begin{equation}\label{C26-two-copy-unit}
                  q(t)q([b](t))\in1+tR[[t]].
\end{equation}
Such a $b$ exists: the Teichm\"uller group
$\mu_{p^2-1}\subset\mathcal O_{K_p}^{\times}$ contains an element
with $b^{p+1}=-1$, because $2(p+1)$ divides $p^2-1$.

For actual nonzero tame translated reduced theta sections, write the
factorization already proved in C18--C19 as
\[
 s_T(t)=q(t)u_T(t),\qquad u_T(t)\in R[[t]]^{\times},
\]
after a finite extension carrying the prime-to-$p$ torsion data.
Then the positive-degree two-section pullback
\begin{equation}\label{C26-theta-pullback}
                 f(t)=s_{T_1}(t)s_{T_2}([b](t))
\end{equation}
is an integral unit whenever $b\bar b=-1$.
Consequently
\begin{equation}\label{C26-strong}
 v_p(D^nf(0))\geq
     \max\left\{0,\left\lfloor\frac{np}{p^2-1}\right\rfloor-1\right\}.
\end{equation}
The same statement holds for positive tensor powers of this
pullback.  It concerns a specified isotropic formal direction;
C22's one-copy derivative has no such cancellation.

\paragraph{The global CM isogeny quotient is exactly rational.}
First take a nonzero $a\in\mathcal O$ such that
$a\equiv1\pmod{2\mathcal O}$ and $p\nmid a\bar a$.
Write $d=a\bar a\in\Z_{>0}$; it is odd.
Let $\Lambda$ be the complex period lattice of the algebraic
differential.  Bannai--Kobayashi, arXiv:math/0610163v4,
Theorem~1.7, Definition~1.8 and Example~1.9, pp.~10--11,
give the exact factors of automorphy
\[
 \theta(z+\gamma)
 =\alpha(\gamma)
   \exp\!\left(\pi H(z,\gamma)+\frac\pi2H(\gamma,\gamma)\right)
   \theta(z),\qquad
 H(z,w)=\frac{z\bar w}{\pi A},
\]
where $\alpha(\gamma)=-1$ off $2\Lambda$ and equals one on
$2\Lambda$.  Since $a\Lambda\subset\Lambda$,
$a\equiv1\pmod2$ gives $a\gamma\equiv\gamma\pmod{2\Lambda}$.
Thus $\alpha(a\gamma)=\alpha(\gamma)=\alpha(\gamma)^d$.
Also $H(az,a\gamma)=dH(z,\gamma)$.
All the factors of automorphy therefore cancel in
\[
                 R_a(z)=\frac{\theta(az)}{\theta(z)^d}.
\]
This is a meromorphic function on the elliptic curve, with
\begin{equation}\label{C26-rational-divisor}
       \operatorname{div}(R_a)=[a]^*(O)-d(O),
       \qquad R_a(z)=a z^{1-d}(1+O(z)).
\end{equation}
The function is algebraic over the field of definition of the
CM endomorphism and the differential.  Indeed the displayed
divisor is defined there.  The vector space of rational sections
with this prescribed principal divisor is one-dimensional after
extension to $\C$; ordinary field-extension compatibility of
global sections descends this one-dimensional space.  Fixing the
nonzero Laurent leading coefficient to be $a$ selects a unique
element over the original field.  This also shows that no
uncontrolled transcendental scalar enters $R_a$.

\paragraph{The primitive local quotient is an integral unit.}
The isogeny $[a]$ on the good model has degree $d$ prime to $p$.
Its kernel is finite \'etale, and its nonidentity geometric
points do not specialize to the identity.  Equivalently,
$[a](t)/(a t)\in R[[t]]^{\times}$ and the only kernel point
in the identity formal disc is $O$.
It follows from \eqref{C26-rational-divisor} that the rational
function
\[
              V_a(t)=\frac{R_a(\lambda(t))t^{d-1}}a
\]
has no horizontal zero or pole on that formal disc, and has
constant term one.  Here the notation $R_a(\lambda(t))$ means
the expansion of the algebraic function $R_a$ in the parameter
$t$.

For precision, the completed local ring is the regular factorial
ring $R[[t]]$.  After cancelling the order $1-d$ at the identity,
the divisor of $V_a$ in this local ring is supported, if anywhere,
on its sole vertical height-one prime $(\varpi_R)$.
Hence $V_a=\varpi_R^m v$ for a unit
$v\in R[[t]]^{\times}$ and some integer $m$.
Evaluation at $t=0$ gives $1=\varpi_R^m v(0)$, forcing $m=0$.
Thus $V_a\in1+tR[[t]]$.
The factoriality used here is the regular-local-ring theorem;
an exact reference is Stacks Project, Tag~0AG0.

The logarithm intertwines CM actions:
$\lambda([a](t))=a\lambda(t)$.  Therefore
\begin{align*}
 \frac{q([a](t))}{q(t)^d}
 &=R_a(\lambda(t))\frac{t^d}{[a](t)}\\
 &=V_a(t)\frac{a t}{[a](t)}\in1+tR[[t]].
\end{align*}
This proves \eqref{C26-unit-isogeny} for the specified global
CM elements.  The argument used the actual canonical reduced
theta and its algebraic Laurent normalization, not just its
divisor class.

\paragraph{Passage to the completed CM action.}
For a fixed $b\in\mathcal O_{K_p}^{\times}$, choose
$a_j\in\mathcal O$ with
\[
 a_j\equiv1\pmod{2\mathcal O},\qquad
 a_j\longrightarrow b\quad\hbox{in }\mathcal O_{K_p}.
\]
The Chinese remainder theorem and maximality at $p$ supply
such a sequence, and all its sufficiently late terms have
$p\nmid a_j\bar a_j$.  C19's completed action gives
coefficientwise convergence $[a_j](t)\to[b](t)$ in $R[[t]]$;
also $a_j\bar a_j\to b\bar b$ in $\Z_p$.

For every fixed power of $t$, the coefficient of $q(t)^u$
is a finite sum of binomial polynomials $\binom{u}{m}$
multiplied by fixed coefficients of $(q-1)^m$.
It is therefore continuous in $u\in\Z_p$, even though the
coefficients of $q$ need not be integral.
Composition with $[a_j](t)$ is likewise coefficientwise
continuous, because these series have zero constant term.
It follows that
\[
 \frac{q([a_j](t))}{q(t)^{a_j\bar a_j}}
          \longrightarrow
 \frac{q([b](t))}{q(t)^{b\bar b}}
\]
coefficientwise.  Every series on the left has integral
coefficients and constant term one.  The valuation ring $R$
is closed in its completion, so the limit lies in
$1+tR[[t]]$.  This proves \eqref{C26-unit-isogeny} and
then \eqref{C26-two-copy-unit}.

\paragraph{Actual theta tensors and their differential contraction.}
The integral-unit factorization of the nonzero tame translates comes
from Katagiri--Kobayashi, arXiv:2609.07024v1,
Definition~3.1 and Remark~3.2, p.~9, and Lemma~3.3(1),
pp.~9--10, together with C19's identification of the reduced
Gaussian factor.  Their constants are units in this normalization.
Using \eqref{C26-two-copy-unit}, the right side of
\eqref{C26-theta-pullback} becomes
\[
       \bigl(q(t)q([b](t))\bigr)
          u_{T_1}(t)u_{T_2}([b](t)),
\]
a product of integral units.  C19 now proves
\eqref{C26-strong}.  Positive integral powers remain integral
units, so the same estimate applies to them.

More generally, take units $b_1,\ldots,b_r$ in
$\mathcal O_{K_p}$ and nonnegative integers $d_1,\ldots,d_r$
satisfying
\[
                  \sum_i d_i b_i\bar b_i=0.
\]
Multiplying \eqref{C26-unit-isogeny} with these powers cancels
the entire remaining exponent of $q$.  Hence
$\prod_i s_{T_i}([b_i](t))^{d_i}$ is an integral unit and
satisfies the same strong derivative estimate.
This is a rank-one isotropic direction with the stated
unit coordinates; no assertion for arbitrary nonunit columns
or higher-rank blocks is used.

The map $t\mapsto([b_1](t),\ldots,[b_r](t))$ is an integral
formal group homomorphism.  Its invariant tangent vector is
$(b_1,\ldots,b_r)$, so its $n$th derivative is precisely the
contraction by
$(\sum_i b_iD_i)^n$ of the product theta section at the origin.
This identifies the differential being estimated.  For
$b_i\notin\mathcal O$, the map is formal and need not
algebraize to a map of complex elliptic curves.
The global CM approximants above justify the finite-jet limit.
\hfill$\square$

\paragraph{Relation to the remaining bridge.}
This supplies an explicit positive-degree canonical-frame
cancellation on an actual CM elliptic formal group, including
its primitive integral normalization.  The isotropic condition
is additional geometric data, so it does not conflict with
C22.  Matching a full archive determinant block, its unit minor
and its normalizing factors still requires the exact dictionary
in A8 and the archive comparison.  Nor does this local
theta identity identify the integral cycle with a
Kudla--Rapoport moduli condition.
